% Exercice: Exercice 1
% Sujet: Analyse
% Généré automatiquement

\begin{exercice}[Analyse]
\label{ex:2}

1.a) Calculer la $limite de$\\
en +0\\
0,2^\{5\} pt\\
\item b) Calculer f'(r) où f' est la fonction dérivée de f .
0,5 pt\\
2.a) Dresser le tableau de variations de\\
[0; +œ[.\\
0,5 pt\\
su\\
\item b) Construire (C,).
0,5 pt\\
$Uo = 1 et pour$\\
\item 3. Soient (U,) et (\%;) les suites numériques définies respeclivement
par\\
3un\\
tout n € IN ,\\
Un+1\\
$V = 1 +$\\
3+4Un\\
5^\{0\}n+3\\
Montrer que pour tout entier naturel n, V + 4 =\\
0,2^\{5\} pt\\
est une suite arithmé$tique de raison 4 et de premier terme Vo = 4$.0,7^{5} pt\\
4.a) Montrer que\\
\item b) $Exprimer Vn en fonction de n pour tout entier naturel$ n.
0,2^\{5\} pt\\
En déduire Un en fonction de n-\\
0,5 pt\\

\end{exercice}

% --- Fin Exercice 1 ---
